\section{State topology and bounded sequence quotients}
\label{sec:preliminaries}

All algebras and Hilbert spaces are over \(\C\).
A MASA is automatically a von Neumann algebra: its relative commutant
equals itself, by taking real and imaginary parts.
Partitions are finite indexed families of orthogonal projections with
sum \(\one\); we allow zero projections.
For a self-adjoint operator \(t\), let \(t_+\) denote its positive part.
We use the same functional-calculus notation for continuous real
functions, including
\[
 \clip_c(t)=\max(-c,\min(t,c)),\qquad c>0.
\]

Suppose for this section that \(M\) has a faithful normal state
\(\varphi\), and \(A\subseteq M_\varphi\) is a MASA in its centralizer.
Here
\[
 M_\varphi=\{a\in M:\varphi(az)=\varphi(za)\text{ for all }z\in M\}.
\]
The modular group fixes \(A\) pointwise. The conditional-expectation
theorem gives a faithful normal \(\varphi\)-preserving expectation
\(E_A:M\to A\); see \cite{Takesaki72} or \cite[Theorem~6.6]{Hiai20}.
Set
\[
 \norm{z}_\varphi=\varphi(z^*z)^{1/2},
 \qquad
 \norm{z}_\varphi^\#=\norm{z}_\varphi+\norm{z^*}_\varphi.
\]
On bounded sets these seminorms determine, respectively, the
\(\sigma\)-strong and \(\sigma\)-strong* topologies.
Indeed, in the faithful state representation,
\(\norm{z}_\varphi=\norm{z\xi_\varphi}\), and \(M'\xi_\varphi\) is
dense. Boundedness therefore promotes convergence on \(\xi_\varphi\)
to strong convergence on all vectors. The bounded
\(\sigma\)-strong topology is intrinsic to the von Neumann algebra.
Apply the same reasoning to adjoints for the second assertion.
In standard \(L^2(M)\) notation we write
\(\xi_\varphi=\varphi^{1/2}\). Elements of \(A\) commute with this
vector under the left and right actions.

\begin{lemma}\label{lem:compression-net}
Every \(A\)-partition compression \(C_{\mathcal Q}\) preserves
\(\varphi\) and is contractive for both state seminorms.
As \(\mathcal Q\) refines through all finite \(A\)-partitions,
\[
 C_{\mathcal Q}(z)\longrightarrow E_A(z)
 \quad\text{in }\norm{\cdot}_\varphi^\#
 \qquad(z\in M).
\]
The same holds along any refining sequence of partitions generating \(A\).
\end{lemma}
\begin{proof}
The centralizer identity gives
\(\varphi(C_{\mathcal Q}(z))=\varphi(z)\).
Compression is unital completely positive, so its Schwarz inequality
gives state-norm contraction; apply it to adjoints as well.
On the state Hilbert space the compression is implemented by the
orthogonal projection
\[
 \zeta\longmapsto\sum_{q\in\mathcal Q}q\zeta q.
\]
These projections decrease under refinement. Their action on
\(z\varphi^{1/2}\) is
\(C_{\mathcal Q}(z)\varphi^{1/2}\).
Meanwhile the operators \(C_{\mathcal Q}(z)\) are uniformly bounded.
Every ultraweak cluster point commutes with all projections of \(A\),
and hence is in \(A\). Since
\(E_A C_{\mathcal Q}(z)=E_A(z)\), the cluster point equals \(E_A(z)\).
The Hilbert-space limit of the decreasing projections must therefore
be \(E_A(z)\varphi^{1/2}\).
Repeating for \(z^*\) proves the claim.
For a generating refining sequence, commuting with all of its
partition projections is already equivalent to commuting with \(A\),
so the same proof applies.
\end{proof}

The following quotient allows us to record vanishing state errors
while retaining norm estimates. Its multiplier condition matters
when \(\varphi\) is not tracial.

\Needspace{12\baselineskip}
\begin{lemma}\label{lem:sequence-quotient}
Let
\[
 \mathcal I_\varphi=
 \{(z_j)\in\ell^\infty(M):\norm{z_j}_\varphi^\#\to0\},
\]
and let
\[
 \mathcal C_\varphi=
 \{t\in\ell^\infty(M):
 t\mathcal I_\varphi\subseteq\mathcal I_\varphi,\
 \mathcal I_\varphi t\subseteq\mathcal I_\varphi\}.
\]
Then \(\mathcal C_\varphi\) is a unital \(C^*\)-algebra and
\(\mathcal I_\varphi\) is a closed two-sided ideal in it.
Constant sequences from \(M\) and all sequences in
\(\ell^\infty(A)\) belong to \(\mathcal C_\varphi\).
\end{lemma}
\begin{proof}
The set \(\mathcal I_\varphi\) is self-adjoint and norm closed.
If \(s,t\in\mathcal I_\varphi\), then
\[
 \norm{s_jt_j}_\varphi\leq
 \norm{s_j}\norm{t_j}_\varphi,\qquad
 \norm{(s_jt_j)^*}_\varphi\leq
 \norm{t_j}\norm{s_j^*}_\varphi ,
\]
so \(st\in\mathcal I_\varphi\).
The two multiplier conditions now show directly that
\(\mathcal C_\varphi\) is a unital, norm-closed \(*\)-algebra,
contains \(\mathcal I_\varphi\), and makes it an ideal.
Multiplication by a fixed bounded operator is continuous on bounded
sets for strong* convergence, so constant sequences are multipliers.

If \(a\in A\), then left multiplication has state-norm bound
\(\norm{a}\), and the centralizer property gives
\[
 \norm{ta}_\varphi^2
 =\varphi(a^*t^*ta)
 =\varphi(t^*t aa^*)
 \leq \norm{a}^2\varphi(t^*t).
\]
The last inequality follows from positivity, since
\(\varphi(t^*t b)=\varphi(b^{1/2}t^*t b^{1/2})\geq0\) for
\(b\geq0\) in \(A\).
The same estimates apply to adjoints, uniformly on bounded
\(A\)-sequences. Thus these sequences are multipliers.
\end{proof}

Write \([t_j]\) for an image in
\(\mathcal C_\varphi/\mathcal I_\varphi\).
A sequence of \(A\)-partitions of uniformly bounded size acts by
compression in this quotient.
If \(q\) refines \(p\), its compression applied to a \(p\)-compression
is its compression of the original element.
Since these maps are unital and positive, a common refinement
preserves separate scalar upper and lower bounds for a self-adjoint
element. In particular it preserves a norm bound.

We will also use the following direct consequence of functional calculus.
If a self-adjoint multiplier sequence \(T_j\) has
\(\norm{[T_j]}\leq c\), then
\[
 w_j=T_j-\clip_c(T_j)\in\mathcal I_\varphi,\qquad
 \norm{T_j-w_j}\leq c.
\]
If \(T_j\) is block diagonal for a partition, so is \(w_j\).
This converts the quotient estimates below into explicit corrections.
